\section{Tool-Selection Performance}\label{sec:exp:performance}

This section supports Claim 1 --- rotational steering improves multi-tool tool selection. Six methods are compared under the same prompt, scorer, and greedy decoding: \texttt{no\_steer}, our stationary-K implementation \texttt{ocq\_bias\_a0.3}, external canonical SEKA, canonical AdaSEKA, and our query-coverage (\texttt{ocq\_qbias\_b-0.03}) and layer-adaptive K+Q (\texttt{ocq\_ladapt\_k0.05\_q-0.03}).

\subsection{MetaTool Subtask4}\label{sec:exp:metatool}

MetaTool Subtask4 (N=497) assigns exactly two ground-truth tools per query, making it the cleanest benchmark for sequential coverage. The locked rows of Table~\ref{tab:e1-subtask4} already tell the central story. On the same catalog-derived ontology basis, the stationary key-side operator shows a strongly negative sign ($-4.57$pp), while our two operators --- on the same basis and rank but with the sign flipped from key amplification to query subtraction --- deliver positive gains: $+2.28$pp (Q-coverage) and $+2.08$pp (layer-adaptive K+Q). External canonical SEKA/AdaSEKA rows are filled after the reproduction gate of Appendix~\ref{app:seka-repro}.

\begin{table}[t]
\centering
\caption{MetaTool Subtask4 (N=497), Qwen2.5-7B-Instruct. \textbf{Bold} = locked; \emph{italic} = placeholder with expected direction.}
\label{tab:e1-subtask4}
\small
\begin{tabularx}{\linewidth}{@{}Xrrr@{}}
\toprule
Method & F1 & Exact & $\Delta$F1 \\
\midrule
\texttt{no\_steer} & \textbf{0.7307} & \textbf{0.5252} & --- \\
canonical SEKA amp=1.0 & \emph{placeholder} & \emph{placeholder} & $\ll 0$ \\
canonical AdaSEKA amp=1.0 & \emph{placeholder} & \emph{placeholder} & $\ll 0$ \\
Stationary K \texttt{a0.3} & \textbf{0.6850} & \textbf{0.4728} & $\mathbf{-4.57}$ \\
Q-coverage (ours) & \textbf{0.7535} & \textbf{0.5332} & $\mathbf{+2.28}$ \\
Layer-adaptive K+Q (ours) & \textbf{0.7514} & \textbf{0.5473} & $\mathbf{+2.08}$ \\
\bottomrule
\end{tabularx}
\end{table}

\subsection{$\tau^2$-bench retail}\label{sec:exp:tau2}

For cross-domain validation we evaluate first-turn action-set selection on $\tau^2$-bench retail, where ground-truth tool counts range over $k\in[0,13]$ in contrast to MetaTool's $k=2$ --- an environment that exposes the regime-split prediction. Even the N=30 smoke in Table~\ref{tab:e2-tau2} shows the expected pattern: Q-coverage raises Qwen F1 by $+3.74$pp while layer-adaptive drops by $-3.73$pp. The asymmetry is not noise but the quantitative signature predicted by Corollary~\ref{cor:regime} --- the marginal benefit of K-imprinting decays on long sequences. The same operator that dominates at $k=2$ weakens at $k\ge 4$; this is what the theory already anticipated. The N=114 full run will provide statistical confirmation.

\begin{table}[t]
\centering
\caption{$\tau^2$-bench retail. Qwen2.5-7B. N=30 smoke is locked.}
\label{tab:e2-tau2}
\small
\begin{tabularx}{\linewidth}{@{}lXrrr@{}}
\toprule
N & Method & F1 & Recall & $\Delta$F1 \\
\midrule
\multirow{5}{*}{30 (smoke)}
 & \texttt{no\_steer} & \textbf{0.5735} & \textbf{0.5340} & --- \\
 & Q-coverage (ours) & \textbf{0.6109} & \textbf{0.5717} & $\mathbf{+3.74}$ \\
 & Layer-adaptive K+Q (ours) & \textbf{0.5362} & \textbf{0.4831} & $\mathbf{-3.73}$ \\
 & canonical SEKA amp=1.0 & \emph{placeholder} & \emph{placeholder} & $\le 0$ \\
 & canonical AdaSEKA amp=1.0 & \emph{placeholder} & \emph{placeholder} & $\le 0$ \\
\midrule
\multirow{3}{*}{114 (full)}
 & \texttt{no\_steer} & \emph{placeholder} & \emph{placeholder} & --- \\
 & Q-coverage (ours) & \emph{placeholder} & \emph{placeholder} & $> 0$ \\
 & Layer-adaptive K+Q (ours) & \emph{placeholder} & \emph{placeholder} & $\approx 0$ or $<0$ \\
\bottomrule
\end{tabularx}
\end{table}

\subsection{Why the gain appears: position-level decomposition}\label{sec:exp:mechanism}

That the aggregate direction matches the prediction is necessary but not explanatory. The real evidence for \emph{why} the gain exists lives in a position-level decomposition. Theorem~\ref{thm:no-memory}'s observable consequence is ``the stationary key-side leads the model to re-select the first tool at the second step,'' so a specific footprint must appear at a specific position. \texttt{repeated\_first\_tool\_rate} must rise under SEKA above \texttt{no\_steer} and fall under our operators, while at the same time \texttt{second\_distinct\_hit\_rate} must rise under our operators. Only if all three observations hold simultaneously does the ``erasing repetition'' interpretation find support in data.

Figure~\ref{fig:stepwise} and Table~\ref{tab:e3-mechanism} show the pattern. Lining four methods against four position-level metrics, both stationary key-side operators (ours and canonical SEKA) raise \texttt{repeated\_first\_tool\_rate} noticeably above \texttt{no\_steer} and show the sharpest drop on \texttt{second\_distinct}. Our operators do the opposite: they weaken repetition while raising second-position accuracy. The \emph{joint directionality} of these four numbers is the central evidence. Had we only had F1, one could argue that ``the same basis happened to produce different numbers''; when four metrics point to the same story, the coincidence argument is no longer available.

\begin{figure}[t]
\centering
\includegraphics[width=0.88\linewidth]{fig3_stepwise.pdf}
\caption{Position-level hit rates. Stationary K-side methods raise \texttt{repeated\_first} above \texttt{no\_steer}; our operators lower it while raising \texttt{second\_distinct} --- the direct confirmation of Corollary~\ref{cor:seka-repeat}.}\label{fig:stepwise}
\end{figure}

\begin{table}[t]
\centering
\caption{Position-level hit rates on MetaTool Subtask4 N=497. Arrows indicate the direction predicted by the theorems.}
\label{tab:e3-mechanism}
\small
\begin{tabularx}{\linewidth}{@{}Xcccc@{}}
\toprule
Method & first\_hit & second\_hit & second\_distinct & repeated\_first $\downarrow$ \\
\midrule
\texttt{no\_steer} & \emph{baseline} & \emph{baseline} & \emph{baseline} & \emph{baseline} \\
canonical SEKA amp=1.0 & $\approx$ & $\downarrow\downarrow$ & $\downarrow\downarrow$ & $\boldsymbol{\uparrow\uparrow}$ \\
Stationary K \texttt{a0.3} & $\approx$ & $\downarrow\downarrow$ & $\downarrow\downarrow$ & $\boldsymbol{\uparrow\uparrow}$ \\
Q-coverage (ours) & $\approx$ & $\uparrow$ & $\uparrow$ & $\downarrow$ \\
Layer-adaptive K+Q (ours) & $\boldsymbol{\uparrow}$ & $\boldsymbol{\uparrow\uparrow}$ & $\boldsymbol{\uparrow\uparrow}$ & $\downarrow$ \\
\bottomrule
\end{tabularx}
\end{table}

\subsection{Ruling out alternative explanations}\label{sec:exp:ablations}

Three alternative explanations can threaten the central claim. First, that the gain is due to basis quality. Second, that the effect is simply because our operator happens to be ``gentler'' than the controls. Third, that our operator acts as a kind of confidence booster where stronger application means better results. Each is ruled out by a specific ablation, and we close with an empirical check of the diagnostic bound in Theorem~\ref{thm:bound}.

\paragraph{Basis specificity.}
Table~\ref{tab:e4-basis} isolates the ``due to the basis'' interpretation. On feature-shuffled and random bases the stationary K operator collapses completely to zero. The same operator retains its ``clean negative effect'' of $-4.57$pp only on the catalog ontology basis, meaning the basis direction fully determines stationary-K behavior. Q-coverage retains F1 above 0.7068 under basis replacement, remaining relatively less sensitive. The PCA-of-K row decides the last step: if PCA is indistinguishable from real, the strong ``catalog ontology is part of the operator'' reading weakens; if PCA is neutral, the claim narrows to ``low-rank structure is necessary but the catalog is optimal.'' In either case the random/shuffled collapse rules out the ``any direction will do'' alternative.

\begin{table}[t]
\centering
\caption{Qwen2.5-7B, Subtask4 N=497: basis $\times$ operator. \textbf{Bold} = locked.}
\label{tab:e4-basis}
\small
\begin{tabularx}{0.9\linewidth}{@{}Xcc@{}}
\toprule
Basis & Q-coverage F1 & Stationary K F1 \\
\midrule
Real $B_{\mathrm{ont}}$ (ours) & \textbf{0.7471} & \textbf{0.6850} \\
Feature-shuffled & \textbf{0.7254} & \textbf{0.0000} \\
Random orthonormal & \textbf{0.7068} & \textbf{0.0000} \\
PCA-of-K (calibration) & \emph{placeholder} & \emph{placeholder} \\
\bottomrule
\end{tabularx}
\end{table}

\paragraph{The ``safer because gentler'' objection.}
A natural pushback is that the real basis beats the null controls because it perturbs less. Table~\ref{tab:e-perturb-mag} rejects this explanation directly. On Qwen L=13 with $\alpha=0.3$ (N=30), the real ontology basis produces a \emph{larger} effective attention-perturbation magnitude than either random or feature-shuffled bases --- mean $621.30$ versus $399.56$ and $291.98$ respectively. Yet the only stable K-side direction is the real one. The real basis wins not because it is weaker but because it is aligned with a specific stable subspace.

\begin{table}[t]
\centering
\caption{Effective attention-perturbation magnitude on Qwen L=13, $\alpha=0.3$, N=30. The real basis perturbs the most yet is the only one that preserves stability.}
\label{tab:e-perturb-mag}
\small
\begin{tabularx}{0.9\linewidth}{@{}Xccc@{}}
\toprule
Basis & Mean & Median & Std \\
\midrule
Real $B_{\mathrm{ont}}$ (ours) & \textbf{621.30} & \textbf{593.71} & 317.59 \\
Random orthonormal & 399.56 & 360.01 & 235.86 \\
Feature-shuffled & 291.98 & 241.87 & 169.97 \\
\bottomrule
\end{tabularx}
\end{table}

\paragraph{The ``more is better'' objection.}
A third pushback is that Q-coverage acts as a confidence booster, so stronger $|\beta|$ should keep improving results. Data rules this out too. Figure~\ref{fig:beta-sweep} shows the narrow operating window of a Q-only $\beta$ sweep on Qwen: a peak near $\beta \approx -0.1$, with $\beta = -0.7$ collapsing F1 completely to zero. The pattern is not ``more bias is better'' but ``a small negative Q-side projection helps and most other settings hurt,'' which matches a coverage interpretation rather than a crude confidence boost.

\begin{figure}[t]
\centering
\includegraphics[width=0.75\linewidth]{fig_qbias_ablations.pdf}
\caption{Q-coverage $\beta$ sweep on Qwen Subtask4. The narrow operating window is inconsistent with the confidence-boost interpretation and consistent with coverage.}\label{fig:beta-sweep}
\end{figure}

\paragraph{Empirical validity of the diagnostic bound.}
Finally we check that the inequality of Theorem~\ref{thm:bound} is not vacuous. On Qwen2.5-7B-Inst (L=13) with the stationary K intervention at $\alpha=0.3$, N=100, and across all 2800 head--query combinations, the measured LHS/RHS ratio has a median of $2.36\times 10^{-8}$ with a 100\% pass rate. The bound is conservative but stable, supporting its use as an informative diagnostic when comparing Q-side and K-side interventions on a common scale.

\paragraph{Behavior in the single-tool regime.}
A final supporting observation reinforces the claim from the opposite direction. On Qwen's MetaTool Subtask1 single-tool regime, a stationary K-bias ($\alpha=0.3$) gives a \emph{large positive gain} of $+11.16$pp. At first this looks inconsistent with our story, but it is the opposite: in a single-tool regime the bias ``repeat the first choice'' is neutral or helpful, which is exactly why the SEKA family excels on single-concept editing. Our two operators are not blanket replacements for stationary K but are designed specifically for the multi-tool problem.

\section{Efficiency}\label{sec:exp:efficiency}

Claim 2 --- that rotational steering is efficient relative to its alternatives --- is supported not by a new experiment but by a direct comparison along four axes. Table~\ref{tab:e-efficiency} places our method against tool-aware fine-tuning, prompt scaffolding, KV-cache optimization, and stationary activation steering.

\begin{table}[t]
\centering
\caption{Efficiency comparison. Our method operates as an inference-time hook with zero training and zero token overhead; storage is limited to the per-head basis, approximately $2.6$ MB on Qwen2.5-7B.}
\label{tab:e-efficiency}
\small
\begin{tabularx}{\linewidth}{@{}Xrrrr@{}}
\toprule
Approach & Training cost & Token overhead & Inference latency & Storage \\
\midrule
Tool-aware fine-tuning (Gorilla, ToolLlama) & $\mathcal{O}(\text{days})$ & 0 & 0 & full model \\
Prompt scaffolding (ReAct) & 0 & $+200\sim 500$ tokens & scales with prompt & 0 \\
KV-cache opt (KIVI) & 0 & 0 & minor decrease & cache $-30\%$ \\
Stationary activation steering (SEKA) & projector training ($\mathcal{O}$(min)) & 0 & minor increase & projector matrix \\
\textbf{This paper (Q-coverage, K+Q)} & \textbf{0} & \textbf{0} & \emph{to measure} & $\boldsymbol{\approx 2.6}$ \textbf{MB} \\
\bottomrule
\end{tabularx}
\end{table}

The ``to measure'' cell is filled by two planned measurements. Inference-latency profiling records per-query wall-clock time under \texttt{no\_steer}, Q-coverage, layer-adaptive K+Q, and canonical SEKA with greedy decoding, and memory profiling verifies that the added GPU footprint matches the theoretical $\approx 2.6$ MB of basis storage. Because the hook adds only an $O(Tdr)$ matmul to each attention call, the expected overhead is under 1\% of total inference time; the measurements convert that prediction into a number. Worth noting is that KV-cache optimization and our method target different objectives (memory versus task performance) and are \emph{composable} --- methods like KIVI compress the cache while our operator improves task-level selection on the same cache.
